%%%%%%%%%%%%%%%%%%%%%%%%%%%%%%%%%%%%%%%%%
% Academic Curriculum Vitae
% XeLaTeX Template
%%%%%%%%%%%%%%%%%%%%%%%%%%%%%%%%%%%%%%%%%

\documentclass[11pt]{article}

\input{structure.tex}

%----------------------------------------------------------------------------------------

\begin{document}

%----------------------------------------------------------------------------------------
%	HEADER
%----------------------------------------------------------------------------------------

\cvheader
    {Haichao Wang}
    {SIGS, Tsinghua University, Shenzhen, China}
    {wanghc23@mails.tsinghua.edu.cn}
    {https://hychaowang.github.io}
    {https://scholar.google.com/citations?user=-exqY3gAAAAJ\&hl=en}

%----------------------------------------------------------------------------------------
%	EDUCATION
%----------------------------------------------------------------------------------------

\section{Education}

\eduentry{2023--2026}
    {M.S.\ in Data Science and Information Technology}
    {\href{https://www.tsinghua.edu.cn/en/}{Tsinghua University}, China}
    {Advisor: Prof.\ Yuxing Han \enspace$\mid$\enspace Research: Efficient Video Computer Vision with Motion Estimation}

\eduentry{2019--2023}
    {B.E.\ in Computer Science and Technology}
    {\href{https://bjtu.edu.cn}{Beijing Jiaotong University}, China}
    {Research: Efficient Physiological Signal Classification and System Design}

%----------------------------------------------------------------------------------------
%	RESEARCH EXPERIENCE
%----------------------------------------------------------------------------------------

\section{Research Experience}

\resentry
    {Physiological Signal Classification and System Design}
    {\href{http://insis.bjtu.edu.cn/}{Institute of Network Science and Intelligent System, BJTU, China}}
    {May 2021 -- Aug. 2023}
    {\begin{itemize}
        \item Proposed a multi-level knowledge distillation framework for lightweight single-channel EEG physiological signal classification.
        \item Proposed a spatial-temporal knowledge distillation framework for lightweightmulti-channel physiological signal classification.
        \item Developed MicroSleepNet, a lightweight model with low computational cost, for a physiological signal classification software.
    \end{itemize}}

\resentry
    {Efficient Video Computer Vision with Motion Estimation}
    {\href{https://www.tsinghua.edu.cn/}{Tsinghua University, China}}
    {Oct. 2023 -- Dec. 2025}
    {\begin{itemize}
        \item Applied video motion estimation to accelerate the video computer vision system.
        \item Proposed an efficient pipeline directly processing Bayer data, removing the ISP latency.
    \end{itemize}}

\resentry
    {Generative Human Motion Mimicking}
    {\href{https://www.cam.ac.uk/}{Tsinghua University, China \& University of Cambridge, UK}}
    {Oct. 2025 -- Present}
    {\begin{itemize}
        \item Developed an algorithm using diffusion trajectory matching for controlled long-range motion generation.
    \end{itemize}}
    

%----------------------------------------------------------------------------------------
%	PUBLICATIONS
%----------------------------------------------------------------------------------------

\section{Publications}

\pubentry
    {Diffusion Trajectory Matching for Controlled Long-Range Motion Generation}
    {\textbf{Haichao Wang}, Alexander Okupnik, Jiangtao Wen, Yuxing Han, Johannes Schneider, Kyriakos Flouris}
    {ECCV, 2026 (in preparation)}

\pubentry
    {Efficient Video Computer Vision with Fast Motion Estimation and Learned Perception Residual}
    {\textbf{Haichao Wang}, Jiangtao Wen, Yuxing Han}
    {ECCV, 2026 (in preparation)}

\pubentry
    {Leveraging motion estimation for Efficient Bayer-Domain Video Convolutional Networks}
    {\textbf{Haichao Wang}, Jiangtao Wen, Yuxing Han}
    {ICIP 2026 (under review)}

\pubentry
    {From Sensor to Semantics: Fast Motion Estimation and Residual Correction for Efficient Bayer-Domain Video Vision}
    {\textbf{Haichao Wang}, Jiangtao Wen, Yuxing Han}
    {IEEE Transactions on Pattern Analysis and Machine Intelligence, 2026 (under review)}



\pubentry
    {Cross-Modal Knowledge Distillation for Enhanced Unimodal Emotion Recognition}
    {Ziyu Jia, Yucheng Liu, \textbf{Haichao Wang}, and Tianzi Jiang}
    {IEEE Transactions on Affective Computing, 2025}

\pubentry
    {Mutual Distillation Extracting Spatial-Temporal Knowledge for Lightweight Multi-Channel Sleep Stage Classification}
    {Ziyu Jia, \textbf{Haichao Wang}, Yucheng Liu, and Tianzi Jiang}
    {SIGKDD, 2024}

\pubentry
    {DistillSleepNet: Heterogeneous Multi-Level Knowledge Distillation via Teacher Assistant for Sleep Staging}
    {Ziyu Jia, Heng Liang, Yucheng Liu, \textbf{Haichao Wang}, and Tianzi Jiang}
    {IEEE Transactions on Big Data, 2024}

\pubentry
    {EmotionKD: A Cross-Modal Knowledge Distillation Framework for Emotion Recognition Based on Physiological Signals}
    {Yucheng Liu, Ziyu Jia, and \textbf{Haichao Wang}}
    {ACM MM, 2023}

\pubentry
    {Teacher Assistant-Based Knowledge Distillation Extracting Multi-Level Features on Single Channel Sleep EEG}
    {Heng Liang$^*$, Yucheng Liu$^*$, \textbf{Haichao Wang}$^*$, and Ziyu Jia {\small ($^*$Equal contribution)}}
    {IJCAI, 2023}
%----------------------------------------------------------------------------------------
%	HONORS AND AWARDS
%----------------------------------------------------------------------------------------

\section{Honors and Awards}

\awardentry{National Prize, National College Students' Innovation and Entrepreneurship Training Program}
\awardentry{Second Prize, Beijing Undergraduate Mathematical Contest in Modeling}
\awardentry{Second Prize, 11th `Challenge Cup' National Undergraduate Curricular Academic Science and Technology Works}

%----------------------------------------------------------------------------------------

\end{document}
